\begin{textsample}{NEG Dim 1 – 2019\_c4 – Score -5.00 – 2019\_c4/enade\_2019\_eng\_flo\_3a2b53498b.txt}  \label{ex:f1_neg_086}
CAPÍTULO 4 PERCEPÇÃO DA PROVA Nas análises feitas neste capítulo, constam as percepções dos concluintes da Área de Engenharia Florestal sobre a prova aplicada no Enade/2019.
Essas percepções foram mensuradas por meio de nove questões que avaliaram desde o grau de dificuldade da prova até o tempo gasto para concluí -la.
As percepções sobre a prova foram relacionadas ao desempenho do estudante, à Grande Região de funcionamento do curso, à Categoria Administrativa e ao tipo de Organização Acadêmica da IES.
O questionário de percepção da prova encontra -se ao final do Anexo VII, onde \textbf{está} a reprodução do exame.
O desempenho dos estudantes foi classificado em quatro quartos.
Para tanto, esse desempenho foi ordenado de forma ascendente.
O percentil 25, P25, também conhecido como primeiro quartil, é a nota de desempenho que deixa um quarto ( 25 \% ) dos valores observados abaixo e três quartos acima.
Na Figura 1, apresenta -se uma ilustração desse conceito.
O quarto inferior de desempenho é composto pelas notas abaixo do primeiro quartil.
Já o percentil 75, P75, também conhecido como terceiro quartil, é o valor para o qual há três quartos ( 75 \% ) dos dados abaixo e um quarto acima dele.
O quarto superior de desempenho é composto pelas notas iguais ou acima do terceiro quartil.
O percentil 50, P50, também conhecido como mediana, é o valor que divide as notas em dois conjuntos de igual tamanho.
O segundo quarto inclui valores entre o primeiro quartil ( P25 ) e a mediana.
O terceiro quarto contém os valores entre a mediana ( P50 ) e o terceiro quartil ( P75 ).
Vale ressaltar que percentis, quartis e medianas são pontos que, não obrigatoriamente, pertencem ao conjunto original de dados, ao passo que os quartos são subconjuntos dos dados originais.
A seguir, serão apresentados gráficos com resultados selecionados, relativos às nove questões avaliadas por grupos de estudantes.
Nas barras dos gráficos, apresentam -se os percentuais de alunos que assinalaram uma das opções ou a soma das porcentagens daqueles que assinalaram duas ( ou três ) delas.
Por exemplo, para as questões 1 e 2, nos gráficos, é apresentada a porcentagem total de participantes que assinalou as opções ( D ) Difícil e ( E ) Muito difícil.
Em \textbf{cada} barra, foram assinalados, também, os extremos do intervalo de confiança de 95 \% como linhas verticais unidas por uma linha horizontal na forma da letra H maiúscula.
O estimador de um parâmetro com um \textbf{certo} nível de confiança ( e.g. 95,0 \% ) deve conter o parâmetro no intervalo de confiança em 95 \% das vezes.
Na comparação entre os estimadores dos parâmetros de duas classes de uma dada categoria ( e.g.
Norte e Nordeste nas Grandes Regiões, de primeiro e último quartos dentro de desempenho, de IES Públicas e Privadas, ou de Universidades e Faculdades ), associados aos seus respectivos intervalos de confiança, diz -se que não há diferença estatisticamente significativa entre os parâmetros das duas categorias quando há interseção entre os intervalos de confiança, mas se os intervalos de confiança são disjuntos, a diferença existe ( para maiores informações, vide Glossário ).
Nas Tabelas do Anexo II, são apresentados os valores absolutos e a distribuição percentual18 das alternativas válidas das nove questões, segundo o mesmo recorte de desempenho dos alunos e da Grande Região de funcionamento do curso, de Categoria Administrativa e do tipo de Organização Acadêmica da IES. 4.1 GRAU DE DIFICULDADE DA PROVA 4.1.1 Componente de Formação Geral Ao avaliarem Qual o grau de dificuldade desta prova na parte de Formação Geral? ( Questão 1 ), 29,1 \% do grupo de inscritos e presentes optaram pelas alternativas Difícil ou Muito difícil.
Entretanto, para quase dois terços dos estudantes ( 63,1 \% ), o Componente de Formação Geral da prova foi considerado com grau de dificuldade Médio ( Gráfico 4.1, Gráfico 4.2 e, no Anexo II, a Tabela II.1 ).
O percentual de estudantes que considerou a prova Difícil ou Muito difícil foi maior na região Sul, atingindo 35,0 \%, enquanto na região Nordeste, o percentual alcançado foi o menor : 23,9 \%.
No Gráfico 4.1, é possível observar que as diferenças entre a região Sul e as regiões Nordeste e Sudeste são estatisticamente significativas.
Nas Grandes Regiões, a 18 Cumpre \textbf{lembrar} uma das convenções para tabelas numéricas ( pág. iii ) sobre a possibilidade de a soma das partes não resultar em 100 \% por questões de arredondamento. proporção de presentes à prova que avaliou o Componente de Formação Geral como de grau de dificuldade Médio \textbf{esteve} entre 57,5 \% na região Sul e 68,3 \% na região Nordeste.
O percentual de alunos que consideraram o componente de Formação Geral da prova Difícil ou Muito difícil foi decrescente em função dos quartos de desempenho : 33,6 \% no primeiro quarto e 26,4 \% no quarto, o grupo de melhor desempenho na prova.
Nos quartos de desempenho intermediários, a proporção de alunos que considerou a prova Difícil ou Muito difícil foi de 32,5 \%, no segundo quarto, e de 23,9 \%, no terceiro quarto.
A diferença entre o terceiro quarto e os quartos inferiores de desempenho são estatisticamente significativas.
Para todos os quartos de desempenho, a alternativa modal para essa pergunta foi a Médio, com valores entre 60,9 \% para os respondentes no primeiro quarto e 68,4 \% para aqueles no terceiro quarto de desempenho.
O percentual de estudantes que considerou o componente de Formação Geral da prova Difícil ou Muito difícil foi maior entre os estudantes de Faculdades, atingindo 45,9 \%, enquanto os de Universidades atingiram a menor incidência : 27,3 \%, havendo diferença estatisticamente significativa entre os dois tipos de organização.
No Gráfico 4.3, é possível observar que existe diferença estatisticamente significativa entre Universidades e Faculdades.
Nas Organizações Acadêmicas, a proporção de presentes à prova que avaliou o Componente de Formação Geral como de grau de dificuldade Médio \textbf{esteve} entre 45,0 \% nos CEFET/IF e 66,7 \% nos Centros Universitários ( ver também a Tabela II.2, no Anexo II ).
O percentual de estudantes que considerou o componente de Formação Geral da prova Difícil ou Muito difícil foi maior entre os estudantes de cursos de IES Privadas ( 43,8 \% ), uma diferença estatisticamente significativa da proporção registrada em IES Públicas ( 27,5 \% ).
Já a proporção de presentes à prova que avaliou o Componente de Formação Geral como de grau de dificuldade Médio foi de 64,4 \% para os de IES Públicas e de 51,1 \% para os de IES Privadas ( ver Gráfico 4.4 ). 4.1.2 Componente de Conhecimento Específico Ao responderem à Questão 2 – Qual o grau de dificuldade desta prova na parte de Componente Específico? – 34,3 \% do grupo de estudantes classificaram -na como Difícil ou Muito difícil.
Além disso, o Componente modal de Conhecimento Específico da prova foi avaliado com grau de dificuldade Médio por 60,3 \% dos alunos ( Gráfico 4.5, Gráfico 4.6 e, no Anexo II, a Tabela II.3 ).
A análise das respostas dos estudantes quanto ao grau de dificuldade do Componente de Conhecimento Específico da prova, agregado por Grande Região, mostra que a diferença entre a maior proporção de alunos que a avaliaram como Difícil ou Muito difícil ocorrida na região Norte ( 37,6 \% ) é estatisticamente significativa em relação à menor proporção, encontrada na região Sudeste ( 27,9 \% ).
Todas as demais diferenças entre as regiões não são estatisticamente significativas.
O percentual de alunos que classificou o grau de dificuldade como Médio, no Componente de Conhecimento Específico, variou de 57,4 \% a 65,7 \% para as regiões Sul e Sudeste, respectivamente.
Considerando -se a avaliação da dificuldade das questões do Componente de Conhecimento Específico da prova, de acordo com o desempenho dos estudantes, observa -se que há diferença estatisticamente significativa dos resultados entre o primeiro e segundo quartos e o quarto superior de desempenho.
As proporções dos que classificaram a parte específica como Difícil ou Muito difícil variaram de 28,8 \% ( último quarto ) a 38,5 \% ( primeiro quarto ).
As demais proporções dessa resposta foram de 38,4 \% e 31,6 \%, nos quartos segundo e terceiro, respectivamente.
Já as proporções dos que responderam que o grau de dificuldade das questões do Componente de Conhecimento Específico da prova foi Médio variaram de 57,2 \% no primeiro quarto a 64,3 \% no último.
Considerando -se a avaliação da dificuldade das questões do Componente de Conhecimento Específico da prova, o percentual de estudantes que as considerou Difícil ou Muito difícil foi maior entre os estudantes de CEFET/IF, alcançando 50,0 \%, enquanto entre os de Universidades foi registrado o menor percentual : 33,2 \%.
Não houve diferença estatisticamente significativa nas proporções entre os tipos de Organização Acadêmica ( ver Gráfico 4.7 ), sendo as questões consideradas Difícil ou Muito difícil por 33,3 \% dos estudantes de Centros Universitários e por 43,4 \% dos de Faculdades.
Nas Organizações Acadêmicas, a proporção de presentes à prova que consideraram o Componente de Conhecimento Específico como sendo de grau de dificuldade Médio \textbf{esteve} entre 47,5 \%, nos CEFET/IF e 66,7 \% nos Centros Universitários ( ver também Tabela II.4, no Anexo II ).
O percentual de estudantes que considerou o Componente Específico da prova Difícil ou Muito difícil foi maior entre os estudantes de cursos de IES Privadas ( 43,3 \% ), uma diferença estatisticamente significativa da proporção registrada em IES Públicas ( 33,3 \% ).
Já a proporção de presentes à prova que classificou este componente como de grau de dificuldade Médio foi de 61,1 \% para os de IES Públicas e de 53,4 \% para os de IES Privadas ( ver Gráfico 4.8 e Tabela II.4, no Anexo II ). 4.2 EXTENSÃO DA PROVA EM RELAÇÃO AO TEMPO TOTAL Indagados quanto à extensão da prova, em relação ao tempo total oferecido para a sua resolução ( Questão 3 ), os estudantes apontaram, com maior incidência, a alternativa que considerava a extensão Adequada, para todas as agregações consideradas ( Gráfico 4.9, Gráfico 4.10, e, no Anexo II, a Tabela II.5 ).
O percentual de alunos que respondeu ser a extensão da prova adequada foi de 68,9 \%.
Já 26,1 \% dos inscritos presentes consideraram que a prova foi Longa ou Muito longa, e 5,1 \% a avaliaram como Curta ou Muito curta.
Dentre as Grandes Regiões, as proporções daqueles que avaliaram a prova como Longa ou Muito longa, em relação ao tempo total destinado à sua resolução, variaram de 23,2 \% na região Norte até 29,6 \% na região Sul.
Não há diferença estatisticamente significativa entre as regiões.
Considerando -se o desempenho dos alunos, nota -se ainda que 68,2 \% consideraram a extensão da prova Adequada, no quarto de desempenho inferior, e 65,5 \%, no de melhor desempenho ( quarto superior ).
Nos quartos intermediários, essa proporção foi de 71,2 \%, no segundo, e de 70,6 \%, no terceiro quarto, proporções decrescentes com o desempenho.
No Gráfico 4.10, pode -se constatar que não há diferença estatisticamente significativa entre as proporções de estudantes que consideraram a prova Longa ou Muito longa em função dos quartos de desempenho, sendo a maior proporção dessa resposta no último quarto de desempenho ( 29,3 \% ), e a menor, no segundo quarto ( 23,6 \% ), não caracterizando uma tendência.
Considerando -se a extensão da prova, o percentual de estudantes que a avaliou como Longa ou Muito longa, em relação ao tempo total destinado à sua resolução, foi maior entre os estudantes de CEFET/IF, alcançando 35,0 \%, enquanto os de Faculdades alcançaram o menor percentual : 22,8 \%, ainda que sem diferença estatisticamente significativa entre todos os tipos de Organização.
Nas Organizações Acadêmicas, os presentes à prova que consideraram a sua extensão Adequada, em relação ao tempo total destinado à sua resolução, atingiram percentuais entre 62,5 \% nos CEFET/IF e 70,6 \% nas Faculdades ( ver também Tabela II.6, no Anexo II ).
O percentual de estudantes que considerou a prova Longa ou Muito longa, em relação ao tempo total destinado à sua resolução, foi maior entre os estudantes de cursos de IES Públicas ( 26,1 \% ), sem uma diferença estatisticamente significativa da proporção registrada em IES Privadas ( 25,3 \% ).
Já a proporção de presentes à prova que considerou a extensão Adequada foi de 68,9 \% para os de IES Públicas e de 69,1 \% para os de IES Privadas ( ver Gráfico 4.12 e Tabela II.6, no Anexo II ). 4.3 COMPREENSÃO DOS ENUNCIADOS DAS QUESTÕES 4.3.1 Componente de Formação Geral Em relação aos enunciados das questões do Componente de Formação Geral ( Questão 4 ), as opiniões foram positivas, já que 74,6 \% dos alunos avaliados consideraram todos ou a maioria dos enunciados das questões claros e objetivos ( Gráfico 4.13, Gráfico 4.14, e, no Anexo II, a Tabela II.7 ).
Na análise regional, as percentagens de estudantes que avaliaram que todos ou a maioria dos enunciados das questões do Componente de Formação Geral \textbf{estavam} claros e objetivos variaram de 70,6 \% na região Sul a 80,0 \%, na região Nordeste.
A análise das percepções dos estudantes sobre a clareza e a objetividade dos enunciados permite afirmar que todos, ou a maioria dos enunciados de questões relativas ao Componente de Formação Geral, foram considerados claros e objetivos para a maior parte dos respondentes ( maior ou igual a 70,6 \%, em todas as regiões, e maior ou igual a 68,8 \% para todos os quartos de desempenho ).
Segundo o desempenho, observa -se que as proporções dos estudantes que emitiram essa opinião crescem conforme o desempenho aumenta, com exceção do segundo para o terceiro quartos, com diferenças estatisticamente significativas do último para os demais quartos de desempenho.
No quarto superior, a clareza e a objetividade de todos ou da maioria dos enunciados das questões foram percebidas por 82,6 \%.
A proporção dos presentes que considerou que todos ou a maioria dos enunciados das questões do Componente de Formação Geral \textbf{estavam} claros e objetivos foi maior entre os estudantes de CEFET/IF, alcançando 85,0 \%, enquanto entre os de Faculdades a proporção atingida foi a menor : 67,6 \%, ainda que sem diferença estatisticamente significativa entre as proporções de todos os tipos de Organização. ( ver também Tabela II.8, no Anexo II ).
O percentual de estudantes que considerou que todos ou a maioria dos enunciados das questões do Componente de Formação Geral \textbf{estavam} claros e objetivos foi maior entre os estudantes de cursos de IES Públicas ( 75,3 \% ), sem diferença estatisticamente significativa da proporção registrada em IES Privadas ( 67,8 \% ).
Um maior \textbf{detalhamento} pode ser encontrado no Gráfico 4.16 e na Tabela II.8, no Anexo II. 4.3.2 Componente de Conhecimento Específico Em relação aos enunciados das questões do Componente de Conhecimento Específico da prova, para 84,9 \% dos estudantes avaliados da Área de Engenharia Florestal, a clareza e a objetividade ( Questão 5 ) \textbf{estavam} presentes em todas ou na maioria das questões ( Gráfico 4.17 a Gráfico 4.20, e no Anexo II, as Tabelas II.9 e II.10 ).
A maioria dos estudantes de todas as Grandes Regiões brasileiras considerou claros e objetivos todos ou a maioria dos enunciados das questões do Componente de Conhecimento Específico da prova, percentual sempre maior ou igual a 80,6 \%.
A diferença entre as regiões não é estatisticamente significativa.
A proporção de estudantes que considerou os enunciados das questões claros e objetivos apresenta uma tendência crescente em relação ao aumento de desempenho : mais elevada no quarto superior ( 92,4 \% ), se comparada ao quarto inferior de desempenho ( 76,2 \% ).
Há diferença estatisticamente significativa entre o primeiro quarto de desempenho e os demais quartos superiores e entre o último quarto de desempenho e os quartos inferiores.
A proporção dos presentes que considerou que todos ou a maioria dos enunciados das questões do Componente de Conhecimento Específico \textbf{estavam} claros e objetivos foi maior entre os estudantes de CEFET/IF, alcançando 87,5 \%, enquanto entre os de Centros Universitários a proporção alcançada foi a menor : 33,3 \%, ainda que sem diferença estatisticamente significativa entre os Centros Universitários e as demais Organizações Acadêmicas.
No Gráfico 4.19, é possível observar que a diferença entre Faculdades ( 75,0 \% ) e Universidades ( 85,7 \% ) é estatisticamente significativa ( ver também Tabela II.10, no Anexo II ).
O percentual de estudantes que considerou que todos ou a maioria dos enunciados das questões do Componente de Conhecimento Específico \textbf{estavam} claros e objetivos foi maior entre os estudantes de cursos de IES Públicas ( 85,9 \% ), com diferença estatisticamente significativa da proporção em IES Privadas ( 75,8 \% ).
Ver Gráfico 4.20 e Tabela II.10, no Anexo II, para maior \textbf{detalhamento}. 4.4 SUFICIÊNCIA DAS INFORMAÇÕES/INSTRUÇÕES FORNECIDAS Ao avaliarem as informações/instruções fornecidas para a resolução das questões ( Questão 6 ), 87,8 \% dos respondentes da Área de Engenharia Florestal de todo o Brasil afirmaram que estas eram até excessivas ou suficientes em todas ou na maioria das questões ( Gráfico 4.21, Gráfico 4.22 e, no Anexo II, a Tabela II.11 ).
Quanto à distribuição de respondentes pelas Grandes Regiões, observa -se que a proporção de estudantes que considerou as informações/instruções fornecidas até excessivas ou suficientes em todas ou na maioria das questões foi sempre superior ou igual a 80,4 \%, chegando a 93,4 \% na região Sudeste.
Existem diferenças estatisticamente significativas entre a região Norte e as demais regiões, com exceção da região Nordeste, e entre a região Nordeste e a Sudeste.
Levando -se em conta o desempenho dos participantes, observa -se uma tendência crescente das proporções de participantes que avaliaram as informações/instruções como até excessivas ou suficientes em todas ou na maioria das questões, havendo diferenças estatisticamente significativas entre alguns quartos de desempenho : do primeiro em relação ao terceiro e ao último quarto ; e entre o segundo e o último quartos de desempenho, como mostrado no Gráfico 4.22.
O percentual foi mais elevado no quarto superior ( 94,1 \% ), percentual superior ao da média nacional ( 87,8 \% ).
Já no quarto inferior, a suficiência das informações/instruções declarada como até excessiva, em todas ou na maioria das questões foi percebida por 79,2 \% dos respondentes.
Considerando -se as informações/instruções fornecidas para a resolução das questões ( Questão 6 ), a proporção que considerou que estas eram até excessivas ou suficientes em todas ou na maioria das questões foi maior entre os estudantes de Centros Universitários, para os quais a proporção foi de 100,0 \%, enquanto entre os de Faculdades a proporção alcançada foi a menor : 80,7 \%, com diferença estatisticamente significativa entre esses dois tipos de Organização Acadêmica.
No Gráfico 4.23, é possível observar que a diferença entre Centros Universitários e Universidades ( 88,2 \% ) também é estatisticamente significativa ( ver também Tabela II.12, no Anexo II ).
O percentual de estudantes que considerou que as informações/instruções fornecidas para a resolução das questões eram até excessivas ou suficientes em todas ou na maioria das questões foi maior entre os estudantes de cursos de IES Públicas ( 88,5 \% ), uma diferença estatisticamente significativa da proporção registrada em IES Privadas ( 80,8 \% ). ( ver Gráfico 4.24 e Tabela II.12, no Anexo II ). 4.5 DIFICULDADE ENCONTRADA AO RESPONDER À PROVA Indagados sobre as dificuldades com as quais se depararam ao responder à prova ( Questão 7 ), 17,5 \% dos estudantes apontaram o Desconhecimento do conteúdo.
Para 51,0 \%, a Forma diferente de abordagem do conteúdo foi indicada como dificuldade.
Já a Falta de motivação para fazer a prova foi a dificuldade apontada por 10,3 \% dos respondentes.
Considerando -se todo o Brasil, 18,6 \% dos respondentes afirmaram que não tiveram qualquer tipo de dificuldade para responder à prova ( Tabelas II.13 e II.14, no Anexo II ).
Nos Gráficos 4.25 a 4.28, são apresentados os percentuais de estudantes que apontaram o Desconhecimento do conteúdo como dificuldade ao responder à prova.
Na análise por Grandes Regiões, os percentuais de inscritos e presentes que apontaram o Desconhecimento do conteúdo como dificuldade ao responder à prova não superaram 25,6 \%.
Os percentuais variaram de 14,0 \% na região Sul a 25,6 \% na região Nordeste.
A Forma diferente de abordagem do conteúdo foi a escolha modal dos estudantes, com percentuais que variaram de 46,9 \% ( região Nordeste ) a 55,2 \% ( região Norte ).
O percentual de alunos que citou a Falta de motivação como dificuldade variou de 9,2 \% ( região Sudeste ) a 12,1 \% ( região Centro-Oeste ).
Os que declararam não ter qualquer tipo de dificuldade para responder à prova variaram de 15,0 \% na região Norte a 21,7 \% na região Centro-Oeste.
Em relação aos quartos de desempenho, o Desconhecimento do conteúdo foi a opção escolhida por 16,4 \% dos estudantes do quarto inferior e por 19,7 \% do quarto superior, caracterizando uma tendência crescente.
Não são observadas diferenças estatisticamente significativas entre os quartos de desempenho.
A alternativa modal para os alunos, quando agregados pelos quartos de desempenho, para a dificuldade encontrada, foi causada pela Forma diferente de abordagem do conteúdo : 51,6 \%, no quarto inferior e 47,3 \% do último quarto escolheram essa alternativa.
Na análise por Organização Acadêmica, os percentuais de inscritos e presentes que apontaram o Desconhecimento do conteúdo como dificuldade ao responder à prova não superaram 27,5 \%.
Os percentuais variaram de 0,0 \%, para Centros Universitários a 27,5 \%, nos CEFET/IF.
A Forma diferente de abordagem do conteúdo foi a escolha modal dos estudantes, com percentuais que variaram de 49,4 \% ( Universidades ) a 66,7 \% ( Centros Universitários e Faculdades ).
O percentual de alunos que citou a Falta de motivação como dificuldade variou de 0,0 \% ( Centros Universitários e CEFET/IF ) a 11,1 \% ( Universidades ).
Os que declararam não ter qualquer tipo de dificuldade para responder à prova variaram de 7,5 \% nos CEFET/IF a 33,3 \% nos Centros Universitários.
No Gráfico 4.27, é possível observar as diferenças estatisticamente significativas entre os Centros Universitários e as demais Organizações Acadêmicas ( ver também Tabela II.14, no Anexo II ).
Em relação às Categorias Administrativas, o Desconhecimento do conteúdo foi a opção escolhida por 18,0 \% dos estudantes nas IES Públicas e por 13,1 \% nas IES Privadas, sem diferenças estatisticamente significativas entre as duas categorias.
A alternativa modal para os alunos, quando agregados pelas Categorias Administrativas, para a dificuldade encontrada, foi causada pela Forma diferente de abordagem do conteúdo : 49,9 \% para as Públicas e 61,9 \% para as Privadas ( ver Gráfico 4.28 e Tabela II.14, no Anexo II ). 4.6 CONTEÚDOS DAS QUESTÕES OBJETIVAS DA PROVA Ao se analisarem os conteúdos das questões objetivas da prova ( Questão 8 ), verifica -se que um percentual muito pequeno dos estudantes avaliados, apenas 1,5 \%, afirmou que não estudou ainda a maioria desses conteúdos ( Gráficos 4.29 a Gráfico 4.32, e as Tabelas II.15 e II.16, no Anexo II ).
A maioria ( 76,7 \% ) afirmou ter estudado e aprendido muitos ou todos os conteúdos.
Na análise por Grande Região, a proporção de respondentes que escolheu a opção não estudou ainda a maioria desses conteúdos foi pequena.
Observa -se que, nas regiões Norte ( 1,8 \% ), Nordeste ( 3,1 \% ) e Centro-Oeste ( 1,8 \% ), as proporções foram maiores do que a proporção da média nacional ( 1,5 \% ).
Nas regiões Sudeste ( 0,5 \% ) e Sul ( 0,9 \% ), a situação foi inversa : registrou -se proporção mais baixa.
As diferenças entre as regiões não são estatisticamente significativas.
Em todas as regiões, a maioria dos presentes afirmou ter estudado e aprendido muitos ou todos os conteúdos, com proporções variando entre 68,1 \% na região Norte a 87,2 \% na região Sudeste.
Considerando -se separadamente as opiniões de estudantes dos quatro quartos de desempenho, observa -se que, no quarto inferior, 4,0 \% ofereceram como resposta que não estudou ainda a maioria desses conteúdos, sendo 0,0 \% os do quarto superior com a mesma resposta.
A diferença entre os percentuais de alunos que optaram por esse motivo de dificuldade no quarto inferior em relação aos dos demais quartos de desempenho é estatisticamente significativa, bem como entre os do quarto superior e os dos demais quartos.
Tendo -se em conta o quarto superior, 89,4 \% dos alunos afirmaram ter estudado e aprendido muitos ou todos os conteúdos.
No outro extremo, no primeiro quarto, 62,5 \% optaram pelas mesmas categorias.
As proporções são crescentes com o desempenho.
Na análise por Organização Acadêmica, a proporção de respondentes escolhendo a opção não estudou ainda a maioria desses conteúdos foi pequena.
Observa -se que, apenas nas Faculdades ( 4,4 \% ), as proporções foram maiores do que a proporção da média nacional ( 1,5 \% ).
Em Universidades ( 1,3 \% ), Centros Universitários ( 0,0 \% ) e CEFET/IF ( 0,0 \% ), a situação foi inversa : a proporção atingida foi mais baixa.
No Gráfico 4.31, é possível observar que as diferenças de Centros Universitários e CEFET/IF em relação a Faculdades e Universidades são estatisticamente significativas ( ver também Tabela II.16, no Anexo II ).
Em todos os tipos de Organização Acadêmica, a maioria dos presentes afirmou ter estudado e aprendido muitos ou todos os conteúdos, com proporções variando entre 50,0 \% em Centros Universitários a 82,5 \% em CEFET/IF.
Em relação às Categorias Administrativas, observa -se que, nas IES Públicas, a resposta não estudou ainda a maioria desses conteúdos, foi apresentada por 1,3 \%.
Já nas IES Privadas, a mesma resposta foi dada por 3,4 \%.
A diferença registrada nas duas categorias não é estatisticamente significativa.
Considerando -se as IES Públicas, 77,3 \% dos alunos afirmaram ter estudado e aprendido muitos ou todos os conteúdos.
Na outra categoria, 71,3 \% optaram pelas mesmas respostas ( ver Gráfico 4.32 e Tabela II.16 no Anexo II ). 4.7 TEMPO GASTO PARA CONCLUIR A PROVA Ao responderem sobre o tempo de conclusão da prova ( Questão 9 ), mais de quatro quintos dos estudantes ( 83,2 \% ) afirmaram ter gasto Entre duas e quatro horas ( Gráfico 4.33 a Gráfico 4.36 e, no Anexo II, as Tabelas II.17 e II.18 ).
Considerando -se as cinco Grandes Regiões brasileiras, os percentuais dos que utilizaram Entre duas e quatro horas para finalizar a prova nas regiões Norte ( 82,9 \% ), Nordeste ( 80,6 \% ) e Sul ( 82,0 \% ) foram inferiores ao percentual nacional.
Nas demais Grandes Regiões, os percentuais alunos que dispensaram Entre duas e quatro horas para concluir a prova alcançaram percentual mais elevado que o do Brasil : 85,5 \% na região Sudeste e 84,3 \% na região Centro-Oeste, como mostrado no Gráfico 4.33.
A diferença entre as regiões não é estatisticamente significativa.
Uma vez analisadas as alternativas escolhidas pelos estudantes que se situam nos diferentes quartos de desempenho, observa -se uma tendência crescente : uma proporção maior de participantes no quarto superior declarou ter gasto Entre duas e quatro horas para concluir a prova quando comparada com a proporção de participantes nos quartos inferiores ( com exceção do terceiro quarto para o último ).
Há diferenças estatisticamente significativas entre o primeiro quarto e os demais quartos de desempenho.
Na análise por Organização Acadêmica, os que utilizaram Entre duas e quatro horas para finalizar a prova nas Universidades alcançaram percentual inferior ao nacional : 83,0 \% e 83,2 \%, respectivamente.
Nos demais tipos de Organização Acadêmica, o percentual de alunos que dispensou Entre duas e quatro horas para concluir a prova foi mais elevado : 100,0 \% nos Centros Universitários ; 90,0 \% nos CEFET/IF e 83,7 \% nas Faculdades, como mostrado no Gráfico 4.35.
A diferença entre Centros Universitários e as demais organizações acadêmicas é estatisticamente significativa.
Em relação às Categorias Administrativas, observa -se que, nas IES Públicas, 83,5 \% declararam ter gasto Entre duas e quatro horas para concluir a prova, sendo 80,8 \% nas IES Privadas com a mesma resposta.
A diferença registrada nas duas categorias não é estatisticamente significativa ( ver Gráfico 4.36 e Tabela II.18, no Anexo II ).

% matched lemmas: cada, certo, detalhamento, estar, lembrar
\end{textsample}
